\section{Results}
\label{sec:results}

\subsection{The corpus}

The 66 transcripts yield 854 closing pairs. Closings are long: the median is 565
characters and the 90th percentile is 2866, which is already a warning about
pattern matching, since a longer message gives a fixed pattern set more chances
to fire on prose that is not a request.

Thirteen closings (1.5\%) are harness error pages rather than messages the agent
chose to write: an API error or a usage-limit notice left as the last assistant
text. These are excluded from every lane. They matter more than their share
suggests: 4 of the 11 user pushes in the whole corpus follow one. The user was
pushing past an infrastructure failure, and a gate that reads those as premature
stopping is blocking a session that has already ended.

\subsection{What each lane fires on}

\begin{table}[h]
\centering
\begin{tabular}{lrr}
\toprule
Lane configuration & Fires on & Share \\
\midrule
Firm patterns only        & 304 & 36.1\% \\
Any pattern (firm + doubt) & 539 & 64.1\% \\
Judge only                 & 569 & 67.7\% \\
Firm patterns or judge (shipped) & 620 & 73.7\% \\
Firm patterns and judge          & 253 & 30.1\% \\
\bottomrule
\end{tabular}
\caption{Firing rates over the 841 non-noise closings.}
\label{tab:fire}
\end{table}

The shipped configuration fires on nearly three closings in four. That number is
the most uncomfortable result here and we report it without softening. It has
two readings and we cannot separate them with this data. Either the agent
genuinely ends most turns with something nameable left open -- which matches the
authors' experience and is why the gate was built -- or the gate is far too
eager and the tolerable rate is much lower.

What we can say is that the two lanes are not redundant. Firm patterns and the
judge agree on only 253 closings, while their union is 620. Each catches a large
set the other misses, which is the intended division of labour: patterns catch
fixed phrasings cheaply, the judge catches the rest.

\subsection{Cost}

The judge answers in 0.31 seconds at the median and 0.66 at the 90th percentile,
running locally on consumer hardware alongside the agent's own work. The
per-pattern confidence split is what keeps this affordable: 235 closings reach
the judge as doubts rather than being decided by a regular expression, and the
304 firm hits never pay for a model call at all.

\subsection{Against the push label}

We evaluate on the 664 ungated closings, of which 6 non-noise closings were
followed by a bare user push.

\begin{table}[h]
\centering
\begin{tabular}{lrr}
\toprule
Lane configuration & Fires on & Pushes caught \\
\midrule
Firm patterns only & 238 & 3 of 6 \\
Any pattern        & 426 & 6 of 6 \\
Judge only         & 478 & 6 of 6 \\
Firm or judge (shipped) & 514 & 6 of 6 \\
Firm and judge     & 202 & 3 of 6 \\
\bottomrule
\end{tabular}
\caption{Recall against the user-push label, ungated closings only.}
\label{tab:recall}
\end{table}

The shipped configuration catches every push. So does the judge alone, and so
does the pattern lane alone once doubts are included. Firm patterns alone catch
half.

This is a weak result and we want to be exact about why. Six positives cannot
distinguish between configurations that all fire on more than a third of the
corpus; a detector that fired on everything would also score 6 of 6. The label
establishes that no lane is systematically blind to the failure it targets, and
nothing more. Precision is the quantity that matters here and this label cannot
measure it, which is the entire reason for the loop in
Section~\ref{sec:loop}.

\subsection{Before and after the gate}

The corpus splits at 2026-08-28, when the gate was installed. Pushes fall from
10 of 664 ungated closings (1.5\%) to 1 of 190 gated ones (0.5\%).

The direction is what the gate is for, and the result is not significant: a
two-sided Fisher exact test gives $p = 0.47$. With a base rate near one percent,
detecting a halving at conventional power would need several thousand closings
per arm. We report it because the effect size is the one worth powering a future
study on, not because this study establishes it.

\subsection{The judge's confidence}

Taking the probability the judge assigns to its own verdict token, the
distributions overlap heavily: median 0.95 for \textsc{stop} ($n = 569$) and
0.87 for \textsc{ok} ($n = 272$). 115 \textsc{stop} verdicts and 92 \textsc{ok}
verdicts fall below 0.75.

At the extremes the signal is real. On hand-checked cases it reads 0.26 to 0.53
on messages that are finished reports and 0.97 to 1.00 on genuine premature
stops. But on the 30 closings that reach the judge carrying doubts alone, it
stops 24, and among the 6 stopped below 0.75 exactly three are finished reports
and three are real cuts. A threshold there buys one correct release and pays for
it with one lost catch.

The confidence is therefore logged and used by nobody. It is informative to a
maintainer reading a specific case and not safe to act on automatically, and we
think the distinction is worth naming: a signal can be genuine, reproducible,
and still not be a decision rule.
